% SPDX-License-Identifier: GPL-3.0-or-later OR CC-BY-SA-4.0
\chapter{App Ops}\label{ch:app-ops} %%##appendices_appops-chapter==title>>


\section{Background}\label{sec:app-ops-background} %%##appendices_appops-background==title>>
%%!!appendices_appops-background<<
\textbf{App Ops} (short hand for \textbf{Application Operations}) are used by Android system (since Android 4.3) to
control application permissions. The user \emph{can} control some permissions, but only the permissions that are considered
dangerous (and Google thinks knowing your phone number isn't a dangerous thing). So, app ops seems to be the one we need
if we want to install apps like Facebook and it's Messenger (the latter literary records everything if you live outside
the EU) and still want \emph{some} privacy and/or security. Although certain features of app ops were available in
Settings and later in hidden settings in older version of Android, it's completely hidden in newer versions of Android
and is continued to be kept hidden. Now, any app with \textbf{android.Manifest.permission.GET\_APP\_OPS\_STATS}
permission can get the app ops information for other applications but this permission is hidden from users and can only
be enabled using ADB or root. Still, the app with this permission cannot grant or revoke permissions (actually mode of
operation) for apps other than itself (with limited capacity, of course). To modify the ops of other app, the app needs
\textbf{android.Manifest.permission.UPDATE\_APP\_OPS\_STATS} permissions which isn't accessible via \texttt{pm} command.
So, you cannot grant it via root or ADB, the permission is only granted to the system apps. There are very few apps who
support disabling permissions via app ops. The best one to my knowledge is
\href{https://github.com/8enet/AppOpsX}{AppOpsX}. The main (visible) difference between my app (AppManager) and this
app is that the latter also provides you the ability to revoke internet permissions (by writing ip tables). One crucial
problem that I faced during the development of the app ops API is the lack of documentation in English language.
%%!!>>


\section{Introduction to App Ops}\label{sec:introduction-to-app-ops} %%##appendices_appops-introduction==title>>

\begin{figure}[ht]
    \centering
    \includegraphics{./images/appops.svg}
    \caption{How app ops work} %%##appendices_appops-howappopswork==title>>
    \label{fig:appops}
\end{figure}
%%!!appendices_appops-introduction<<
Figure \hyperref[fig:appops]{1} describes the process of changing and processing permission.
\hyperref[sec:appopsmanager]{AppOpsManager} can be used to manage permissions in Settings app. \textbf{AppOpsManager} is
also useful in determining if a certain permission (or operation) is granted to the application. Most of the methods of
\textbf{AppOpsManager} are accessible to the user app but unlike a system app, it can only be used to check permissions
for any app or for the app itself and start or terminating certain operations. Moreover, not all operations are actually
accessible from this Java class. \textbf{AppOpsManager} holds all the necessary constants such as
\hyperref[subsec:op-constants]{\texttt{OP\_*}}, \texttt{OPSTR\_*}, \hyperref[subsec:mode-constants]{\texttt{MODE\_*}}
which describes operation code, operation string and mode of operations respectively. It also holds necessary data
structures such as \hyperref[subsec:package-ops]{PackageOps} and \textbf{OpEntry}. \textbf{PackageOps} holds
\textbf{OpEntry} for a package, and \textbf{OpEntry}, as the name suggests, describes each operation.

\texttt{AppOpService} is completely hidden from a user application but accessible to the system applications.
As it can be seen in Figure \hyperref[fig:appops]{1}, this is the class that does the actual management stuff.
It contains data structures such as \textbf{Ops} to store basic package info and \textbf{Op} which is similar to
\textbf{OpEntry} of \textbf{AppOpsManager}.
It also has \textbf{Shell} which is actually the source code of the \hyperref[sec:appops-cli]{\textit{appops} command line tool}.
It writes configurations to or read configurations from \hyperref[sec:appops-xml]{\texttt{/data/system/appops.xml}}.
System services calls \textbf{AppOpsService} to find out what an application is allowed and what is not allowed to perform,
and \textbf{AppOpsService} determines these permissions by parsing \texttt{/data/system/appops.xml}. If no custom values
are present in \textit{appops.xml}, it returns the default mode available in \textbf{AppOpsManager}.
%%!!>>


\section{AppOpsManager}\label{sec:appopsmanager} %%##appendices_appops-appopsmanager==title>>
%%!!appendices_appops-appopsmanager<<
\href{https://android.googlesource.com/platform/frameworks/base/+/master/core/java/android/app/AppOpsManager.java}{AppOpsManager}
stands for application operations manager. It consists of various constants and classes to modify app operations.

\seealsoinline{\href{https://developer.android.com/reference/android/app/AppOpsManager}{AppOpsManager documentation}}
%%!!>>

\subsection{\texttt{OP\_*} Constants}\label{subsec:op-constants} %%##appendices_appops-op-constants==title>>
%%!!appendices_appops-op-constants_1<<
\texttt{OP\_*} are the integer constants starting from \texttt{0}. \texttt{OP\_NONE} implies that no operations are
specified whereas \texttt{\_NUM\_OP} denotes the number of operations defined in \texttt{OP\_*} prefix. While they
denote each operation, the operations are not necessarily unique. In fact, there are many operations that are actually
a single operation denoted by multiple \texttt{OP\_*} constant (possibly for future use). Vendors may define their own
op based on their requirements. MIUI is one of the vendors who are known to do that.
%%!!>>

\begin{listing}[ht]
    \begin{minted}[frame=lines]{java}
public static final int OP_NONE = -1;
public static final int OP_COARSE_LOCATION = 0;
public static final int OP_FINE_LOCATION = 1;
public static final int OP_GPS = 2;
public static final int OP_VIBRATE = 3;
...
public static final int OP_READ_DEVICE_IDENTIFIERS = 89;
public static final int OP_ACCESS_MEDIA_LOCATION = 90;
public static final int OP_ACTIVATE_PLATFORM_VPN = 91;
public static final int _NUM_OP = 92;
    \end{minted}
    \caption{Sneak-peek of \texttt{OP\_*}} %%##appendices_appops-sneakpeekopconstants==title>>
    \label{lst:op-constants}
\end{listing}

%%!!appendices_appops-op-constants_2<<
Whether an operation is unique is defined by \texttt{sOpToSwitch}.
It maps each operation to another operation or to itself (if it's a unique operation).
For instance, \texttt{OP\_FINE\_LOCATION} and \texttt{OP\_GPS} are mapped to \texttt{OP\_COARSE\_LOCATION}.

Each operation has a private name which are described by \texttt{sOpNames}.
These names are usually the same names as the constants without the \texttt{OP\_} prefix.
Some operations have public names as well which are described by \texttt{sOpToString}.
For instance, \texttt{OP\_COARSE\_LOCATION} has the public name \textbf{android:coarse\_location}.

As a gradual process of moving permissions to app ops, there are already many permissions that are defined under some operations.
These permissions are mapped in \texttt{sOpPerms}.
For example, the permission \textbf{android.Manifest.permission.ACCESS\_COARSE\_LOCATION} is mapped to
\texttt{OP\_COARSE\_LOCATION}. Some operations may not have any associated permissions which have \texttt{null} values.

As described in the previous section, operations that are configured for an app are stored at
\hyperref[sec:appops-xml]{\texttt{/data/system/appops.xml}}.
If an operation is not configured, then whether system will allow that operation is determined from \texttt{sOpDefaultMode}.
It lists the \textit{default mode} for each operation.
%%!!>>

\subsection{\texttt{MODE\_*} Constants}\label{subsec:mode-constants} %%##appendices_appops-mode-constants==title>>
%%!!appendices_appops-mode-constants<<
\texttt{MODE\_*} constants also integer constants starting from \texttt{0}. These constants are assigned to each
operation describing whether an app is authorised to perform that operation. These modes usually have associated names
such as \textbf{allow} for \texttt{MODE\_ALLOWED}, \textbf{ignore} for \texttt{MODE\_IGNORED}, \textbf{deny} for
\texttt{MODE\_ERRORED} (a rather misnomer), \textbf{default} for \texttt{MODE\_DEFAULT} and \textbf{foreground}
for \texttt{MODE\_FOREGROUND}.

\begin{enumerate}
    \addtocounter{enumi}{-1}
    \item \textbf{\texttt{MODE\_ALLOWED}.} The app is allowed to perform the given operation
    \item \textbf{\texttt{MODE\_IGNORED}.} The app is not allowed to perform the given operation, and any attempt to
    perform the operation should \emph{silently fail}, i.e.\ it should not cause the app to crash
    \item \textbf{\texttt{MODE\_ERRORED}.} The app is not allowed to perform the given operation, and this attempt
    should cause it to have a fatal error, typically a \texttt{SecurityException}
    \item \textbf{\texttt{MODE\_DEFAULT}.} The app should use its default security check, specified in \texttt{AppOpsManager}
    \item \textbf{\texttt{MODE\_FOREGROUND}.} Special mode that means ``allow only when app is in foreground.''
    This mode was added in Android 10
    \item \textbf{\texttt{MODE\_ASK}.} This is a custom mode used by MIUI whose uses are unknown.
\end{enumerate}
%%!!>>

\subsection{PackageOps}\label{subsec:package-ops} %%##appendices_appops-packageops==title>>
%%!!appendices_appops-packageops_1<<
\textbf{AppOpsManager.PackageOps} is a data structure to store all the \textbf{OpEntry} for a package. In simple terms,
it stores all the customised operations for a package.
%%!!>>

\begin{listing}[H]
    \begin{minted}[frame=lines]{java}
public static class PackageOps implements Parcelable {
    private final String mPackageName;
    private final int mUid;
    private final List<OpEntry> mEntries;
    ...
}
    \end{minted}
    \caption{Class \texttt{PackageOps}} %%##appendices_appops-classpackageops==title>>
    \label{lst:package-ops-class}
\end{listing}

%%!!appendices_appops-packageops_2<<
As can be seen in Listing \hyperref[cb12]{2}, it stores all \textbf{OpEntry} for a package as well as the
corresponding package name and its kernel user ID\@.
%%!!>>

\subsection{OpEntry}\label{subsec:opentry} %%##appendices_appops-opentry==title>>
%%!!appendices_appops-opentry_1<<
\textbf{AppOpsManager.OpEntry} is a data structure that stores a single operation for any package.
%%!!>>

\begin{listing}[H]
    \begin{minted}[frame=lines]{java}
public static final class OpEntry implements Parcelable {
    private final int mOp;
    private final boolean mRunning;
    private final @Mode int mMode;
    private final @Nullable LongSparseLongArray mAccessTimes;
    private final @Nullable LongSparseLongArray mRejectTimes;
    private final @Nullable LongSparseLongArray mDurations;
    private final @Nullable LongSparseLongArray mProxyUids;
    private final @Nullable LongSparseArray<String> mProxyPackageNames;
    ...
}
    \end{minted}
    \caption{Class \texttt{OpEntry}} %%##appendices_appops-classopentry==title>>
    \label{lst:class-op-entry}
\end{listing}
%%!!appendices_appops-opentry_2<<
Here:
\begin{itemize}
    \item \texttt{mOp}: Denotes one of the \hyperref[subsec:op-constants]{\texttt{OP\_*} constants}
    \item \texttt{mRunning}: Whether the operation is in progress (i.e.\ the operation has started but not finished
    yet). Not all operations can be started or finished this way
    \item \texttt{mMOde}: One of the \hyperref[subsec:mode-constants]{\texttt{MODE\_*} constants}
    \item \texttt{mAccessTimes}: Stores all the available access times
    \item \texttt{mRejectTimes}: Stores all the available reject times
    \item \texttt{mDurations}: All available access durations, checking this with \texttt{mRunning} will tell you for
    how long the app is performing a certain app operation
    \item \texttt{mProxyUids}: No documentation found
    \item \texttt{mProxyPackageNames:} No documentation found
\end{itemize}
%%!!>>

\subsection{Usage}\label{subsec:usage} %%##appendices_appops-usage==title>>
TODO


\section{AppOpsService}\label{sec:appopsservice} %%##appendices_appops-appopsservice==title>>
TODO


\section{appops.xml}\label{sec:appops-xml} %%##appendices_appops-appopsxml==title>>
%%!!appendices_appops-appopsxml_1<<
Latest \texttt{appops.xml} has the following format: (This DTD is made by me and by no means perfect, has compatibility issues.)
%%!!>>

\begin{Verbatim}
<!DOCTYPE app-ops [

<!ELEMENT app-ops (uid|pkg)*>
<!ATTLIST app-ops v CDATA #IMPLIED>

<!ELEMENT uid (op)*>
<!ATTLIST uid n CDATA #REQUIRED>

<!ELEMENT pkg (uid)*>
<!ATTLIST pkg n CDATA #REQUIRED>

<!ELEMENT uid (op)*>
<!ATTLIST uid
n CDATA #REQUIRED
p CDATA #IMPLIED>

<!ELEMENT op (st)*>
<!ATTLIST op
n CDATA #REQUIRED
m CDATA #REQUIRED>

<!ELEMENT st EMPTY>
<!ATTLIST st
n CDATA #REQUIRED
t CDATA #IMPLIED
r CDATA #IMPLIED
d CDATA #IMPLIED
pp CDATA #IMPLIED
pu CDATA #IMPLIED>

]>
\end{Verbatim}
%%!!appendices_appops-appopsxml_2<<
The instruction below follows the exact order given above:
\begin{itemize}
    \item \texttt{app-ops}: The root element. It can contain any number of \texttt{pkg} or package \texttt{uid}
    \begin{itemize}
        \item \texttt{v}: (optional, integer) The version number (default: \texttt{NO\_VERSION} or \texttt{-1})
    \end{itemize}

    \item \texttt{pkg}: Stores package info. It can contain any number of \texttt{uid}
    \begin{itemize}
        \item \texttt{n}: (required, string) Name of the package
    \end{itemize}

    \item Package \texttt{uid}: Stores package or packages info
    \begin{itemize}
        \item \texttt{n}: (required, integer) The user ID
    \end{itemize}

    \item \texttt{uid}: The package user ID. It can contain any number of \texttt{op}
    \begin{itemize}
        \item \texttt{n}: (required, integer) The user ID
        \item \texttt{p}: (optional, boolean) Is the app is a private/system app
    \end{itemize}

    \item \texttt{op}: The operation, can contain \texttt{st} or nothing at all
    \begin{itemize}
        \item \texttt{n}: (required, integer) The op name in integer, i.e.\ AppOpsManager.OP\_*
        \item \texttt{m}: (required, integer) The op mode, i.e.\ AppOpsManager.MODE\_*
    \end{itemize}

    \item \texttt{st}: State of operation: whether the operation is accessed, rejected or running (not available on old versions)
    \begin{itemize}
        \item \texttt{n}: (required, long) Key containing flags and uid
        \item \texttt{t}: (optional, long) Access time (default: \texttt{0})
        \item \texttt{r}: (optional, long) Reject time (default: \texttt{0})
        \item \texttt{d}: (optional, long) Access duration (default: \texttt{0})
        \item \texttt{pp}: (optional, string) Proxy package name
        \item \texttt{pu}: (optional, integer) Proxy package uid
    \end{itemize}
\end{itemize}

This definition can be found at
\href{https://android.googlesource.com/platform/frameworks/base/+/master/services/core/java/com/android/server/appop/AppOpsService.java}{AppOpsService}.
%%!!>>


\section{Command Line Interface}\label{sec:appops-cli}%%##appendices_appops-appopscli==title>>
%%!!appendices_appops-appopscli<<
\texttt{appops} or \texttt{cmd appops} (on latest versions) can be accessible via ADB or root. This is an easier method
to get or update any operation for a package (provided the package name is known). The help page of this command is
self-explanatory:
%%!!>>
\begin{Verbatim}
AppOps service (appops) commands:
help
Print this help text.
start [--user <USER_ID>] <PACKAGE | UID> <OP>
Starts a given operation for a particular application.
stop [--user <USER_ID>] <PACKAGE | UID> <OP>
Stops a given operation for a particular application.
set [--user <USER_ID>] <[--uid] PACKAGE | UID> <OP> <MODE>
Set the mode for a particular application and operation.
get [--user <USER_ID>] <PACKAGE | UID> [<OP>]
Return the mode for a particular application and optional operation.
query-op [--user <USER_ID>] <OP> [<MODE>]
Print all packages that currently have the given op in the given mode.
reset [--user <USER_ID>] [<PACKAGE>]
Reset the given application or all applications to default modes.
write-settings
Immediately write pending changes to storage.
read-settings
Read the last written settings, replacing current state in RAM.
options:
<PACKAGE> an Android package name or its UID if prefixed by --uid
<OP>      an AppOps operation.
<MODE>    one of allow, ignore, deny, or default
<USER_ID> the user id under which the package is installed. If --user is not
specified, the current user is assumed.
\end{Verbatim}
